本研究の主張は、量子計算がVRP/CVRP/EVRPで古典手法より実用的に優位であるというものではない。
むしろ、既存研究がどの技術段階にあり、将来のEV・物流システムが要求する社会段階との間にどのようなギャップが残るかを明示することに意義がある。

本研究の中心的な貢献は、図~\ref{fig:method-workflow}で示した手続きである。
すなわち、個別文献の幅・深さや実行環境を抽出し、それらを量子側準備度の証拠として分類し、社会側SRLと対応づけ、最後に実用候補化に必要な証拠条件へ変換する方法である。
この手続きによって、量子VRP/CVRP/EVRP研究を単なる性能比較ではなく、社会側要求に対する技術準備度のギャップとして評価できる。

幅と深さは、この評価における重要な中間証拠である。
ただし、深さは論文ごとに定義が異なるため、単純比較は避ける必要がある。
また、幅が小さい符号化が常に実行しやすいわけではなく、深さ、ゲート数、実行可能性、ハイブリッドワークフローへの負担を同時に確認する必要がある。

社会側についても、EV保有台数や物流需要を量子計算の直接効果として扱わない。
これらは社会段階を形成する外生・古典側変数であり、量子最適化が影響し得るとしても、経路品質、充電待ち、配送遅延、再最適化頻度などを通じた間接効果に限られる。

本文の表~\ref{tab:transition-summary}と付録表~\ref{tab:transition-results}の状態遷移評価は、この解釈をさらに明確にする。
技術主導シナリオでは、多くのギャップが「条件付き」または「技術側先行」に残り、量子側の証拠だけを1段階改善しても実用候補には届かない。
資源共進化シナリオでも同様に、資源証拠が改善しても、実行可能性、ワークフロー、社会側準備度の判定条件が残る。
これは、量子最適化の実用候補性が量子ビット数や資源推定だけでは決まらないことを示している。

一方、社会需要主導シナリオでは、社会側要求だけが上昇するため、全ギャップが「社会側先行」または「QRL不足」として残る。
これは、EV・物流システム側の要求が高度化しても、量子側のワークフロー、実行可能性、資源証拠が対応しなければ、実用候補性はむしろ遠ざかることを意味する。
したがって、本研究でいう準備度ギャップは、技術不足だけではなく、社会側要求と量子側証拠の時間差・対応不足でもある。

最も重要なのは、技術社会の均衡改善シナリオでのみ実用候補が現れる点である。
ただし、この結果は「量子最適化が実用化可能である」という性能主張ではない。
むしろ、資源、実行可能性、ワークフロー、社会側準備度が同時に改善される場合に限り、実用候補として検討できる条件が形成される、という条件付きの結果である。
この点が、本研究を単なる量子技術ロードマップ比較ではなく、社会側準備度と量子側準備度の共同評価として位置づける理由である。

本研究の方法論的な意義は、ギャップを分類して終わることではなく、次のエンジニアリング上の行動へ接続できる点にある。
例えば、資源ギャップが残る場合には、より具体的な問題インスタンスに対するresource estimationが次の行動になる。
実行可能性ギャップが残る場合には、制約を保つ符号化、feasibility-preserving mixer、後処理、古典ベースラインとの比較が次の検証対象になる。
ワークフローギャップが残る場合には、ショット数、実行時間、最適化反復、TMS・配車データとの接続を含むPoC設計が必要になる。
したがって、SRL--QRLギャップ評価は、実用性を断定するための指標ではなく、PoCや資源推定へ進むための問題分解として機能する。

この実用性は、アンケートやワークショップによって検証できる。
具体的には、表~\ref{tab:gap-matrix}や表~\ref{tab:transition-summary}を提示し、参加者が「どのギャップを先に検証すべきか」「どの証拠を追加すればPoCに進めるか」「社会側期待と技術側証拠のずれが理解できるか」を評価する。
この検証により、本方法論が単なる分析表ではなく、研究計画や実証計画を支援する道具として有効かを確認できる。

本研究には限界もある。
順序尺度スコアと閾値は文献から抽出した証拠に基づくが、実運用データ、TMS・配車ワークフロー、古典ベースライン、費用対効果、実際の交通・物流シミュレーションを直接含むものではない。
また、QRLやSRLは本研究内の分析ラベルであり、NASA TRLやHerrmann et al.のARLと同一ではない。
したがって、本文の表~\ref{tab:transition-summary}は将来の実験・実証の優先順位を整理するための要約表であり、付録表~\ref{tab:transition-results}はその詳細な判定根拠を確認するための表である。
いずれも、量子優位や運用効果を直接示す表ではない。
